% !TeX root = ../main.tex
\subsection{Simulations on two benchmark signals}
\label{sec:simulations}

We compare Bayesian-PMALA with Bayesian and optimization-based quantile trend filtering on two benchmark signals: the varying-smoothness signal (VS) of \citet{onizuka2024fast}, which adds a localized bump to gradual variation, and the smooth sinusoid (S) of \citet{brantley2020baseline}. Section~\ref{sec:sim:individual} considers interior quantile levels, and Section~\ref{sec:sim:extreme} the joint estimation of extreme upper quantiles. Additional tables are in Appendix~\ref{app:simulation_results}.

For $u_i=i/n$, $i=1,\ldots,n$, and $n=500$, the observations are generated as $y_i=f(u_i)+\varepsilon_i$ with
\begin{align}
 f_{\mathrm{VS}}(u)&=2+\sin(4u-2)+2\exp\{-30(4u-2)^2\},
 \label{eq:sim_vs}\\
 f_{\mathrm{S}}(u)&=\sin(2\pi u).
 \label{eq:sim_s}
\end{align}
The errors follow the three noise families of \citet{onizuka2024fast}. The first is Gaussian, $\varepsilon_i\sim\mathcal{N}(0,s(u_i)^2)$ with $s(u)=(1+u^2)/4$. The second is Beta, $\varepsilon_i\sim\mathrm{Beta}(1,11-10u_i)$. The third is mixed normal, $\varepsilon_i\sim u_i\mathcal{N}(-0.2,0.5)+(1-u_i)\mathcal{N}(0.2,0.5)$, with the component variance as second argument. The target is the true conditional quantile $Q_\tau(u)=f(u)+F_u^{-1}(\tau)$, where $F_u$ is the error distribution at $u$.

\subsubsection{Interior quantiles}
\label{sec:sim:individual}

The competitors are the Laplace and horseshoe Gibbs samplers of \citet{onizuka2024fast} and the ADMM estimator of \citet{brantley2020baseline}. The levels are $\tau=0.05$, $0.25$, $0.50$, $0.75$, and $0.95$, and all methods are fitted to the same data in each replication. Bayesian-PMALA uses $\sigma\sim\mathrm{IG}(0.1,0.1)$. For the VS signal, $\lambda\sim\mathrm{Gamma}(1,10^{-3})$, the interweaving step is used, and the chain has $30000$ sweeps of which the first $10000$ are discarded; for the S signal, $\lambda\sim\mathrm{Gamma}(1,5\times10^{-4})$, no interweaving step is used, and the chain has $20000$ sweeps of which the first $5000$ are discarded. The samplers of \citet{onizuka2024fast} use $8000$ sweeps with $2000$ burn-in and an $\mathrm{IG}(1,1)$ prior on $\sigma^2$. The Bayesian point estimates are posterior means. ADMM selects $\lambda$ by extended BIC over a $25$-point logarithmic grid, and a replication is discarded when any candidate on the grid fails to converge.

For each fit we compute $\mathrm{MSE}=n^{-1}\sum_i(\widehat\theta_{\tau,i}-Q_\tau(u_i))^2$, $\mathrm{MAD}=n^{-1}\sum_i|\widehat\theta_{\tau,i}-Q_\tau(u_i)|$, and, for the samplers, the mean width MCIW of the pointwise $95\%$ intervals. All three are averaged over valid replications.

\label{sec:sim:ind:vs}Table~\ref{tab:vs_int_mse} reports the MSE for the VS signal; MAD and MCIW are in Appendix~\ref{app:sim:individual:vs}. In every configuration the MSE values of the three Bayesian samplers are within a factor of $1.22$ of one another, and the horseshoe has the smallest at all five levels under Beta noise. A single global penalty therefore gives nearly the same error as local shrinkage on this signal. ADMM has the largest MSE in $14$ of the $15$ configurations; the exception is mixed-normal noise at $\tau=0.95$. This is consistent with extended BIC selecting a penalty too large for the localized bump, whereas the samplers estimate a smaller one.

\begin{table}[!htbp]
\centering
\small
\caption{Mean squared error (MSE) for the varying-smoothness (VS) signal ($k=2$), $n=500$, at interior quantile levels.}
\label{tab:vs_int_mse}
\begin{tabular}{lrrrrr}
\toprule
Beta & $0.05$ & $0.25$ & $0.50$ & $0.75$ & $0.95$ \\
\midrule
Onizuka (Laplace) & 0.000513 & 0.001024 & 0.001942 & 0.002994 & 0.007269 \\
Onizuka (Horseshoe) & 0.000512 & 0.000943 & 0.001802 & 0.002778 & 0.006493 \\
Bayesian-PMALA & 0.000541 & 0.001032 & 0.001966 & 0.003035 & 0.007911 \\
ADMM & 0.008055 & 0.004120 & 0.006632 & 0.007613 & 0.014649 \\
\midrule
Mixed normal & $0.05$ & $0.25$ & $0.50$ & $0.75$ & $0.95$ \\
\midrule
Onizuka (Laplace) & 0.075391 & 0.031684 & 0.027303 & 0.031800 & 0.077075 \\
Onizuka (Horseshoe) & 0.073529 & 0.031743 & 0.028090 & 0.033608 & 0.079668 \\
Bayesian-PMALA & 0.081173 & 0.033267 & 0.026919 & 0.031086 & 0.074781 \\
ADMM & 0.096462 & 0.092891 & 0.071678 & 0.049944 & 0.068415 \\
\midrule
Gaussian & $0.05$ & $0.25$ & $0.50$ & $0.75$ & $0.95$ \\
\midrule
Onizuka (Laplace) & 0.020916 & 0.008284 & 0.006932 & 0.007558 & 0.018128 \\
Onizuka (Horseshoe) & 0.019990 & 0.008067 & 0.006651 & 0.007358 & 0.017762 \\
Bayesian-PMALA & 0.019703 & 0.007962 & 0.006803 & 0.007469 & 0.018675 \\
ADMM & 0.088108 & 0.025091 & 0.020181 & 0.014800 & 0.022334 \\
\bottomrule
\end{tabular}
\end{table}

\label{sec:sim:ind:smooth}Table~\ref{tab:s_int_mse} reports the MSE for the S signal; MAD and MCIW are in Appendix~\ref{app:sim:individual:smooth}. ADMM has the smallest MSE in every configuration; the MSE of Bayesian-PMALA is between $1.4$ and $3.1$ times as large. This is the smooth-baseline setting for which ADMM was developed. Bayesian-PMALA has a smaller MSE than both samplers of \citet{onizuka2024fast} at $\tau=0.25$, $0.50$, and $0.75$ under every noise family, by $15\%$ to $23\%$ relative to the better of the two. At $\tau=0.05$ and $0.95$ its MSE is larger than both in five of the six configurations, by $8\%$ to $18\%$ relative to the better of the two; the exception is Beta noise at $\tau=0.05$. This signal has no localized feature to which local shrinkage scales could adapt.

\begin{table}[!htbp]
\centering
\small
\caption{Mean squared error (MSE) for the smooth (S) signal ($k=2$), $n=500$, at interior quantile levels.}
\label{tab:s_int_mse}
\begin{tabular}{lrrrrr}
\toprule
Beta & $0.05$ & $0.25$ & $0.50$ & $0.75$ & $0.95$ \\
\midrule
Onizuka (Laplace) & 0.000295 & 0.000794 & 0.001592 & 0.002311 & 0.005225 \\
Onizuka (Horseshoe) & 0.000423 & 0.000812 & 0.001554 & 0.002270 & 0.005342 \\
Bayesian-PMALA & 0.000251 & 0.000649 & 0.001255 & 0.001832 & 0.006169 \\
ADMM & 0.000183 & 0.000425 & 0.000789 & 0.001078 & 0.002464 \\
\midrule
Mixed normal & $0.05$ & $0.25$ & $0.50$ & $0.75$ & $0.95$ \\
\midrule
Onizuka (Laplace) & 0.071456 & 0.026312 & 0.021974 & 0.027079 & 0.067077 \\
Onizuka (Horseshoe) & 0.074248 & 0.026449 & 0.021783 & 0.027401 & 0.072335 \\
Bayesian-PMALA & 0.079543 & 0.022323 & 0.017629 & 0.021790 & 0.072843 \\
ADMM & 0.030277 & 0.010993 & 0.008228 & 0.010769 & 0.023342 \\
\midrule
Gaussian & $0.05$ & $0.25$ & $0.50$ & $0.75$ & $0.95$ \\
\midrule
Onizuka (Laplace) & 0.014933 & 0.006463 & 0.005581 & 0.006136 & 0.014457 \\
Onizuka (Horseshoe) & 0.016242 & 0.006414 & 0.005364 & 0.006003 & 0.014982 \\
Bayesian-PMALA & 0.016324 & 0.005068 & 0.004114 & 0.004793 & 0.015674 \\
ADMM & 0.006561 & 0.003129 & 0.002456 & 0.002843 & 0.005929 \\
\bottomrule
\end{tabular}
\end{table}

\subsubsection{Joint estimation of extreme upper quantiles}
\label{sec:sim:extreme}

The joint sampler of Section~\ref{subsec:joint} fits $\tau\in\{0.90,0.93,0.96,0.97,0.98,0.99\}$ together. The individual sampler, the two samplers of \citet{onizuka2024fast}, and ADMM are fitted one level at a time to the same data. Tables~\ref{tab:vs_ext_upper_mse} and \ref{tab:s_ext_upper_mse} report the MSE at all six levels, and Appendix~\ref{app:sim:joint_extreme} counts crossings among the four levels from $0.96$ to $0.99$.

On the VS signal the joint sampler has the smallest MSE under Gaussian and mixed-normal noise at every level from $0.93$ upward, and its advantage grows as $\tau$ approaches one: at $\tau=0.99$ under mixed-normal noise its MSE is $0.1240$, against $0.2123$ for the horseshoe sampler and $0.2980$ for ADMM. At $\tau=0.90$ under these two noise families the samplers of \citet{onizuka2024fast} have the smaller MSE. Under Beta noise the horseshoe sampler has the smallest MSE from $\tau=0.90$ to $0.98$, as at interior levels, and at $\tau=0.99$ the joint sampler and the individual sampler have the smallest MSE, both at $0.0132$.

On the S signal ADMM has the smallest MSE under Beta and Gaussian noise, with the joint sampler second. Under mixed-normal noise the ADMM error grows from $0.0158$ at $\tau=0.90$ to $0.2911$ at $\tau=0.99$, whereas that of the joint sampler grows from $0.0280$ to $0.0940$; the joint sampler has the smallest MSE from $\tau=0.97$ upward.

On both signals the joint sampler has a smaller MSE than the individual sampler in every configuration but two, both for the VS signal: with Gaussian noise at $\tau=0.90$ the two values are $0.0135$ and $0.0134$, and with Beta noise at $\tau=0.99$ both are $0.0132$. The posterior means of the joint sampler never cross, whereas every individually fitted method produces crossing curves at these closely spaced levels.

\begin{table}[!htbp]
\centering
\small
\caption{Mean squared error at extreme upper quantiles for the VS signal, $n=500$ and $k=2$. The joint sampler fits the six levels together; the other methods fit each level separately.}
\label{tab:vs_ext_upper_mse}
\begin{tabular}{lrrrrrr}
\toprule
Beta & $0.90$ & $0.93$ & $0.96$ & $0.97$ & $0.98$ & $0.99$ \\
\midrule
Onizuka (Laplace) & 0.0048 & 0.0058 & 0.0087 & 0.0102 & 0.0120 & 0.0156 \\
Onizuka (Horseshoe) & 0.0044 & 0.0054 & 0.0075 & 0.0088 & 0.0104 & 0.0138 \\
Bayesian-PMALA & 0.0055 & 0.0068 & 0.0106 & 0.0124 & 0.0154 & 0.0132 \\
ADMM & 0.0174 & 0.0160 & 0.0169 & 0.0174 & 0.0199 & 0.0183 \\
Joint sampler (non-crossing) & 0.0052 & 0.0057 & 0.0083 & 0.0113 & 0.0126 & 0.0132 \\
\midrule
Mixed normal & $0.90$ & $0.93$ & $0.96$ & $0.97$ & $0.98$ & $0.99$ \\
\midrule
Onizuka (Laplace) & 0.0445 & 0.0562 & 0.0992 & 0.1313 & 0.1692 & 0.2437 \\
Onizuka (Horseshoe) & 0.0497 & 0.0627 & 0.0940 & 0.1183 & 0.1441 & 0.2123 \\
Bayesian-PMALA & 0.0541 & 0.0716 & 0.1203 & 0.1468 & 0.1840 & 0.2458 \\
ADMM & 0.0508 & 0.0570 & 0.0926 & 0.1244 & 0.1845 & 0.2980 \\
Joint sampler (non-crossing) & 0.0520 & 0.0548 & 0.0760 & 0.0880 & 0.1034 & 0.1240 \\
\midrule
Gaussian & $0.90$ & $0.93$ & $0.96$ & $0.97$ & $0.98$ & $0.99$ \\
\midrule
Onizuka (Laplace) & 0.0116 & 0.0144 & 0.0225 & 0.0285 & 0.0350 & 0.0500 \\
Onizuka (Horseshoe) & 0.0116 & 0.0142 & 0.0212 & 0.0259 & 0.0318 & 0.0453 \\
Bayesian-PMALA & 0.0134 & 0.0167 & 0.0271 & 0.0328 & 0.0410 & 0.0502 \\
ADMM & 0.0211 & 0.0224 & 0.0234 & 0.0265 & 0.0321 & 0.0407 \\
Joint sampler (non-crossing) & 0.0135 & 0.0141 & 0.0188 & 0.0223 & 0.0265 & 0.0300 \\
\bottomrule
\end{tabular}
\end{table}

\begin{table}[!htbp]
\centering
\small
\caption{Mean squared error at extreme upper quantiles for the S signal, $n=500$ and $k=2$. The joint sampler fits the six levels together; the other methods fit each level separately.}
\label{tab:s_ext_upper_mse}
\begin{tabular}{lrrrrrr}
\toprule
Beta & $0.90$ & $0.93$ & $0.96$ & $0.97$ & $0.98$ & $0.99$ \\
\midrule
Onizuka (Laplace) & 0.0036 & 0.0043 & 0.0061 & 0.0079 & 0.0097 & 0.0122 \\
Onizuka (Horseshoe) & 0.0037 & 0.0045 & 0.0062 & 0.0075 & 0.0091 & 0.0119 \\
Bayesian-PMALA & 0.0042 & 0.0054 & 0.0088 & 0.0107 & 0.0134 & 0.0121 \\
ADMM & 0.0022 & 0.0025 & 0.0032 & 0.0033 & 0.0039 & 0.0050 \\
Joint sampler (non-crossing) & 0.0029 & 0.0030 & 0.0038 & 0.0041 & 0.0049 & 0.0068 \\
\midrule
Mixed normal & $0.90$ & $0.93$ & $0.96$ & $0.97$ & $0.98$ & $0.99$ \\
\midrule
Onizuka (Laplace) & 0.0392 & 0.0491 & 0.0906 & 0.1254 & 0.1657 & 0.2409 \\
Onizuka (Horseshoe) & 0.0433 & 0.0554 & 0.0860 & 0.1102 & 0.1389 & 0.2003 \\
Bayesian-PMALA & 0.0485 & 0.0629 & 0.1119 & 0.1420 & 0.1765 & 0.2324 \\
ADMM & 0.0158 & 0.0182 & 0.0346 & 0.0714 & 0.1732 & 0.2911 \\
Joint sampler (non-crossing) & 0.0280 & 0.0301 & 0.0430 & 0.0512 & 0.0643 & 0.0940 \\
\midrule
Gaussian & $0.90$ & $0.93$ & $0.96$ & $0.97$ & $0.98$ & $0.99$ \\
\midrule
Onizuka (Laplace) & 0.0092 & 0.0112 & 0.0177 & 0.0239 & 0.0317 & 0.0435 \\
Onizuka (Horseshoe) & 0.0096 & 0.0119 & 0.0178 & 0.0229 & 0.0285 & 0.0406 \\
Bayesian-PMALA & 0.0110 & 0.0143 & 0.0238 & 0.0293 & 0.0378 & 0.0460 \\
ADMM & 0.0045 & 0.0054 & 0.0071 & 0.0086 & 0.0111 & 0.0216 \\
Joint sampler (non-crossing) & 0.0072 & 0.0079 & 0.0107 & 0.0121 & 0.0150 & 0.0219 \\
\bottomrule
\end{tabular}
\end{table}
